\documentclass[acmsmall]{acmart}

% workaround for duplicate symbols:
\let\oldBbbk\Bbbk%
\let\Bbbk\undefined%

% lhs2TeX:
%include polycode.fmt

%%
%% \BibTeX command to typeset BibTeX logo in the docs
\AtBeginDocument{%
  \providecommand\BibTeX{{%
    Bib\TeX}}}

%% Rights management information.  This information is sent to you
%% when you complete the rights form.  These commands have SAMPLE
%% values in them; it is your responsibility as an author to replace
%% the commands and values with those provided to you when you
%% complete the rights form.
\setcopyright{acmlicensed}
\copyrightyear{2018}
\acmYear{2018}
\acmDOI{XXXXXXX.XXXXXXX}

%%
%% These commands are for a JOURNAL article.
\acmJournal{JACM}
\acmVolume{37}
\acmNumber{4}
\acmArticle{111}
\acmMonth{8}

%%
%% Submission ID.
%% Use this when submitting an article to a sponsored event. You'll
%% receive a unique submission ID from the organizers
%% of the event, and this ID should be used as the parameter to this command.
%%\acmSubmissionID{123-A56-BU3}

%%
%% For managing citations, it is recommended to use bibliography
%% files in BibTeX format.
%%
%% You can then either use BibTeX with the ACM-Reference-Format style,
%% or BibLaTeX with the acmnumeric or acmauthoryear sytles, that include
%% support for advanced citation of software artefact from the
%% biblatex-software package, also separately available on CTAN.
%%
%% Look at the sample-*-biblatex.tex files for templates showcasing
%% the biblatex styles.
%%

%%
%% The majority of ACM publications use numbered citations and
%% references.  The command \citestyle{authoryear} switches to the
%% "author year" style.
%%
%% If you are preparing content for an event
%% sponsored by ACM SIGGRAPH, you must use the "author year" style of
%% citations and references.
%% Uncommenting
%% the next command will enable that style.
\citestyle{acmauthoryear}

% Preamble:
${ for(header-includes) }
${ header-includes }
${ endfor }

%%
%% end of the preamble, start of the body of the document source.
\begin{document}

%%
%% The "title" command has an optional parameter,
%% allowing the author to define a "short title" to be used in page headers.
% Document title
${ if(title) }
${ if(shorttitle) }
\title[${ shorttitle }]{${ title }}
${ else }
\title{${ title }}
${ endif }
${ else }
\title{\textcolor{red}{TODO: Set title.}}
${ endif }

%%
%% The "author" command and its associated commands are used to define
%% the authors and their affiliations.
%% Of note is the shared affiliation of the first two authors, and the
%% "authornote" and "authornotemark" commands
%% used to denote shared contribution to the research.
${ for(author) }
${ if(it.name) }
\author{${ it.name }}
${ else }
\author{\textcolor{red}{TODO: Author name.}}
${ endif }
${ if(it.note) }
\authornote{${ it.note }}
${ endif }
${ if(it.email) }
\email{${ it.email }}
${ endif }
${ if(it.orchid) }
\orcid{${ it.orchid }}
${ endif }
\affiliation{%
  ${ if(it.institute) }
  \institution{${ it.institute }}
  ${ endif }
  ${ if(it.city) }
  \city{${ it.city }}
  ${ endif }
  ${ if(it.state) }
  \state{${ it.state }}
  ${ endif }
  ${ if(it.country) }
  \country{${ it.country }}
  ${ endif }
}
${ endfor }

%%
%% By default, the full list of authors will be used in the page
%% headers. Often, this list is too long, and will overlap
%% other information printed in the page headers. This command allows
%% the author to define a more concise list
%% of authors' names for this purpose.
${ if(shortauthors) }
\renewcommand{\shortauthors}{${ shortauthors}}
${ endif }

%%
%% The abstract is a short summary of the work to be presented in the
%% article.
${ if(abstract) }
\begin{abstract}
  ${ abstract }
\end{abstract}
${ endif }

%%
%% The code below is generated by the tool at http://dl.acm.org/ccs.cfm.
%% Please copy and paste the code instead of the example below.
%%
\begin{CCSXML}
<ccs2012>
 <concept>
  <concept_id>00000000.0000000.0000000</concept_id>
  <concept_desc>Do Not Use This Code, Generate the Correct Terms for Your Paper</concept_desc>
  <concept_significance>500</concept_significance>
 </concept>
 <concept>
  <concept_id>00000000.00000000.00000000</concept_id>
  <concept_desc>Do Not Use This Code, Generate the Correct Terms for Your Paper</concept_desc>
  <concept_significance>300</concept_significance>
 </concept>
 <concept>
  <concept_id>00000000.00000000.00000000</concept_id>
  <concept_desc>Do Not Use This Code, Generate the Correct Terms for Your Paper</concept_desc>
  <concept_significance>100</concept_significance>
 </concept>
 <concept>
  <concept_id>00000000.00000000.00000000</concept_id>
  <concept_desc>Do Not Use This Code, Generate the Correct Terms for Your Paper</concept_desc>
  <concept_significance>100</concept_significance>
 </concept>
</ccs2012>
\end{CCSXML}

\ccsdesc[500]{Do Not Use This Code~Generate the Correct Terms for Your Paper}
\ccsdesc[300]{Do Not Use This Code~Generate the Correct Terms for Your Paper}
\ccsdesc{Do Not Use This Code~Generate the Correct Terms for Your Paper}
\ccsdesc[100]{Do Not Use This Code~Generate the Correct Terms for Your Paper}

%%
%% Keywords. The author(s) should pick words that accurately describe
%% the work being presented. Separate the keywords with commas.
% Mandatory.
${ if(keywords) }
\keywords{${ keywords }}
${ else }
\keywords{\textcolor{red}{TODO: Set keywords}}
${ endif }

\received{20 February 2007}
\received[revised]{12 March 2009}
\received[accepted]{5 June 2009}

%%
%% This command processes the author and affiliation and title
%% information and builds the first part of the formatted document.
\maketitle

${ body }

%%
%% The acknowledgments section is defined using the "acks" environment
%% (and NOT an unnumbered section). This ensures the proper
%% identification of the section in the article metadata, and the
%% consistent spelling of the heading.
% Optional. Acknowledgements.
${ if(acknowledgements) }
\begin{acks}
${ acknowledgements }
\end{acks}
${ endif }

%%
%% The next two lines define the bibliography style to be used, and
%% the bibliography file.
\bibliographystyle{ACM-Reference-Format}
${ if(bibliography) }
\bibliography{${ bibliography }}
${ endif }

% TODO: Support appendices.

\end{document}
\endinput
%%
%% End of file `sample-acmsmall.tex'.
